%! Equivalent Representations of Polynomial and Rational Expressions — Unit 1, Lesson 1.11 — Guided Notes (Answer Key)
\documentclass[10pt]{article}
\usepackage{precalculus-article}
\usepackage{precalculus-key}   % pulls in -boxes; do NOT also load -boxes
\newcommand{\TallMath}[1]{$\displaystyle #1\rule[-1.4em]{0pt}{3.2em}$}

\begin{document}

\pageheader{Unit 1, Lesson 1.11}{Guided Notes --- Answer Key}
\namedateperiod

\vspace{0.08in}

\begin{objectivebox}
By the end of this lesson, I will be able to\ldots
\begin{enumerate}[itemsep=2pt]
  \item Rewrite a polynomial or rational expression in \ans{factored} form or \ans{standard} form
        and explain what each form reveals. \hfill\textit{(LO 1.11.A)}
  \item Perform polynomial \ans{long division} to write $f(x) = g(x)q(x) + r(x)$ and use the result
        to find \ans{slant} asymptotes. \hfill\textit{(LO 1.11.B)}
  \item Apply the \ans{binomial} theorem and Pascal's Triangle to expand $(a+b)^n$. \hfill\textit{(LO 1.11.C)}
\end{enumerate}
\end{objectivebox}

\vspace{0.06in}

\begin{vocabbox}
\termblanklong{factored form}
\ansline{A polynomial or rational expression written as a product of linear and/or irreducible quadratic factors.}
\termblanklong{standard form}
\ansline{A polynomial written with terms in descending order of degree; reveals degree, leading coefficient, and end behavior.}
\termblanklong{polynomial long division}
\ansline{An algorithm for dividing polynomial $f$ by polynomial $g$ to find quotient $q$ and remainder $r$ such that $f = gq + r$.}
\termblanklong{quotient}
\ansline{The polynomial $q(x)$ resulting from dividing $f(x)$ by $g(x)$ in long division.}
\termblanklong{remainder}
\ansline{The polynomial $r(x)$ left over after long division, with $\deg(r) < \deg(g)$.}
\termblanklong{slant asymptote}
\ansline{A linear asymptote $y = q(x)$ that occurs when the degree of the numerator is exactly one more than the degree of the denominator.}
\termblanklong{binomial theorem}
\ansline{The formula $(a+b)^n = \sum_{k=0}^{n}\binom{n}{k}a^{n-k}b^k$, which uses Pascal's Triangle coefficients.}
\termblanklong{Pascal's Triangle}
\ansline{A triangular array of numbers where each entry is the sum of the two entries above it; row $n$ gives the coefficients for $(a+b)^n$.}
\end{vocabbox}

\vspace{0.06in}

\begin{hookbox}
Three rational functions are written on the board:
\[
  f(x) = \frac{2x^2 - 8}{x^2 + x - 6}
  \qquad
  g(x) = \frac{2(x-2)(x+2)}{(x-2)(x+3)}
  \qquad
  h(x) = \frac{2(x+2)}{x+3}, \quad x \ne 2
\]
These are the \textit{same function} in three different faces. \ans{Factored} form reveals zeros and
holes. \ans{Standard} form reveals end behavior. The simplified form reveals the \ans{domain restriction} of the
restricted domain.
\end{hookbox}

\vspace{0.06in}

\begin{notesbox}{LO 1.11.A --- Factored Form vs.\ Standard Form}
\renewcommand{\arraystretch}{1.5}
\begin{tabularx}{\linewidth}{@{} >{\bfseries}p{0.22\linewidth} X X @{}}
 & \textbf{Factored Form} & \textbf{Standard Form} \\
\midrule
What it reveals & \ans{zeros, holes, asymptotes, domain} & \ans{end behavior, degree, leading coefficient} \\[4pt]
                & zeros, holes, asymptotes & end behavior, degree, leading coefficient \\[4pt]
Example (poly)  & $(x-1)(x+3)(x-2)$ & $x^3 - 7x + 6$ (expand) \\[4pt]
Example (rational) & $\dfrac{(x-2)(x+2)}{(x-2)(x+3)}$ & $\dfrac{x^2-4}{x^2+x-6}$ \\
\end{tabularx}

\medskip
\textbf{Key idea (EK 1.11.A.1):} Factored form reveals \ans{zeros}, \ans{asymptotes},
\ans{holes}, and domain restrictions.

\textbf{Key idea (EK 1.11.A.2):} Standard form reveals \ans{end behavior and degree}.
\end{notesbox}

\vspace{0.06in}

\begin{notesbox}{LO 1.11.B --- Polynomial Long Division}
\textbf{Division Algorithm:} If $f$ is divided by $g$ (with $\deg g \geq 1$), then
\[
  f(x) = g(x)\cdot q(x) + r(x), \quad \text{where } \deg(r) < \deg(g).
\]

\textbf{Steps:}
\begin{enumerate}[itemsep=3pt]
  \item Write numerator and denominator in \ans{descending} order (fill gaps with $0$ coefficients).
  \item Divide the \ans{leading} term of the dividend by the leading term of the divisor. Write this as the first term of $q(x)$.
  \item Multiply that term by the \ans{entire divisor} and subtract from the dividend.
  \item Bring down the next term and repeat until $\deg(\text{remainder}) < \deg(\text{divisor})$.
\end{enumerate}

\medskip
\textbf{Example:} Divide $f(x) = 2x^2 + 3x - 5$ by $g(x) = x - 1$.

\smallskip
\begin{tabular}{r|l}
\multicolumn{1}{r}{} & \ans{$2x + 5$} \\[2pt]
\cline{2-2}
$x - 1$ & $2x^2 + 3x - 5$ \\
         & $- ($\ans{$2x^2 - 2x$}$)$ \\
\cline{2-2}
         & \ans{$5x - 5$} \\
         & $- ($\ans{$5x - 5$}$)$ \\
\cline{2-2}
         & \ans{$0$} \\
\end{tabular}

\smallskip
Result: $q(x) = $ \ans{$2x+5$}, $r(x) = $ \ans{$0$}, so $f(x) = (x-1)($\ans{$2x+5$}$) + $\ans{$0$}.

\medskip
\textbf{Slant Asymptote (EK 1.11.B.2):} When $\deg(f) = \deg(g) + 1$, the slant asymptote of
$\dfrac{f(x)}{g(x)}$ is $y = q(x)$ (the \ans{quotient polynomial} after long division).
\end{notesbox}

\vspace{0.06in}

\begin{notesbox}{LO 1.11.C --- Binomial Theorem and Pascal's Triangle}
\textbf{Pascal's Triangle} (fill in the missing entries):

\begin{center}
\begin{tabular}{ccccccccccc}
  &   &   &   & 1 &   &   &   &   \\
  &   &   & 1 &   & 1 &   &   &   \\
  &   & 1 &   & 2 &   & 1 &   &   \\
  & 1 &   & 3 &   & 3 &   & 1 &   \\
1 &   & \ans{4} &   & \ans{6} &   & \ans{4} &   & 1 \\
\end{tabular}
\end{center}

\textbf{Rule:} Each entry equals the \ans{sum} of the two entries above it.

\medskip
\textbf{Binomial Theorem (EK 1.11.C.1):}
\[
  (a + b)^n = \sum_{k=0}^{n} \binom{n}{k} a^{n-k} b^k
\]
The coefficients $\binom{n}{k}$ match row $n$ of Pascal's Triangle.

\medskip
\textbf{Example:} Expand $(x + 2)^3$.

Using row 3 of Pascal's Triangle ($1,\ 3,\ 3,\ 1$):
\[
  (x+2)^3 = 1\cdot x^3 + 3\cdot x^2(\ans{$2$}) + 3\cdot x(\ans{$4$}) + 1\cdot(\ans{$8$})
           = \ans{$x^3 + 6x^2 + 12x + 8$}
\]
\end{notesbox}

\end{document}
